\documentclass{article}
\usepackage[utf8]{inputenc}
\usepackage[T1]{fontenc}
\usepackage{geometry}
\geometry{
    a4paper,
    left=2.5cm,
    right=2.5cm,
    top=2.5cm,
    bottom=2.5cm
}
\usepackage{graphicx}
\usepackage{float}
\usepackage{hyperref}
\usepackage[english]{babel}
\usepackage{amsmath}
\usepackage{amssymb}
\usepackage{array}
\usepackage{booktabs}
\usepackage{caption}
\usepackage{subcaption}
\usepackage{enumitem}
\usepackage{fancyhdr}
\usepackage{titlesec}
\usepackage{url}
\usepackage{textcomp}
\usepackage{xcolor}
\usepackage{listings}
\usepackage{multirow}
\usepackage{longtable}
\usepackage{tabularx}
\usepackage{times}

% Set up page style
\pagestyle{fancy}
\fancyhf{}
\rfoot{\thepage}
\renewcommand{\headrulewidth}{0pt}

% Title formatting
\titleformat{\chapter}[display]{\normalfont\bfseries}{\chaptertitle\thechapter}{0pt}{\Huge}
\titlespacing*{\chapter}{0pt}{50pt}{40pt}

% Document info
\title{Template Document}
\author{Academic Support Agent}
\date{\today}

\begin{document}
\maketitle

%
% paper title
% can use linebreaks \\ within to get better formatting as desired
\title{\LaTeX\ Template for Papers to Be Submitted to EuCAP 2016 Conference (\em{Paper Title})}



% author names and affiliations
% use a single column layout for the different affiliations.
% In case the affiliation is too long then use \\ to create multiple lines, as shown in affiliation 4
% Too long lines can result in margin upload warnings ....
\author{\IEEEauthorblockN{
Author One\IEEEauthorrefmark{1},   % 1st author, 1st affiliations
Author Two\IEEEauthorrefmark{2},   % 2nd author, 2nd affiliations
Author Three\IEEEauthorrefmark{3},    % 3rd author, 3rd affiliations
Author Four\IEEEauthorrefmark{4}      % 4th author, 4th affiliations
}                                     % ...
%\\
\IEEEauthorblockA{\IEEEauthorrefmark{1}% 1st affiliations
(Affiliation): dept. name of organization, name of organization, acronyms acceptable, City, Country, e-mail address optional}
\IEEEauthorblockA{\IEEEauthorrefmark{2}% 2nd affiliations
(Affiliation): dept. name of organization, name of organization, acronyms acceptable, City, Country, e-mail address optional}
\IEEEauthorblockA{\IEEEauthorrefmark{3}% 3rd affiliations
(Affiliation): dept. name of organization, name of organization, acronyms acceptable, City, Country, e-mail address optional}
\IEEEauthorblockA{\IEEEauthorrefmark{4}% 4th affiliations
(Affiliation): dept. name of organization, name of organization, acronyms acceptable, City, Country, \\
 e-mail address if desired or not}
}

% conference papers do not typically use \thanks and this command
% is locked out in conference mode. If really needed, such as for
% the acknowledgment of grants, issue a \IEEEoverridecommandlockouts
% after \documentclass



% use for special paper notices
%\IEEEspecialpapernotice{(Invited Paper)}




% make the title area
\maketitle


\begin{abstract}
%\boldmath
This electronic document is a  live  template and already defines the components of your paper [title, text, heads, etc.] in its style sheet. \em{*CRITICAL:  Do Not Use Symbols, Special Characters, or Math in Paper Title or Abstract}. 
\end{abstract}
% IEEEtran.cls defaults to using nonbold math in the Abstract.
% This preserves the distinction between vectors and scalars. However,
% if the conference you are submitting to favors bold math in the abstract,
% then you can use LaTeX's standard command \boldmath at the very start
% of the abstract to achieve this. Many IEEE journals/conferences frown on
% math in the abstract anyway.

% no keywords
{\smallskip \keywords antenna, propagation, measurement.}



% For peer review papers, you can put extra information on the cover
% page as needed:
% \ifCLASSOPTIONpeerreview
% \begin{center} \bfseries EDICS Category: 3-BBND \end{center}
% \fi
%
% For peerreview papers, this IEEEtran command inserts a page break and
% creates the second title. It will be ignored for other modes.
\IEEEpeerreviewmaketitle

\vspace{7pt}
\section{Introduction}
% no \IEEEPARstart
This demo file is intended to serve as a ``starter file''
for EuCAP 2016 conference papers produced under \LaTeX\ using
IEEEtran.cls version 1.7 and later.
% You must have at least 2 lines in the paragraph with the drop letter
% (should never be an issue)
Note that authors must first submit electronically either a short 2 A4 pages version or a full 4 or 5 A4 pages (including figures) version of the paper in Adobe Acrobat PDF format. The maximum size of the full paper is 5 pages. For more details see http://www.eucap2016.org. We wish you the best of success.

\vspace{7pt}
\section{General Guidelines}
Define abbreviations and acronyms the first time they are used in the text, even after they have been defined in the abstract. Abbreviations such as IEEE, SI, MKS, CGS, sc, dc, and rms do not have to be defined. Do not use abbreviations in the title or heads unless they are unavoidable.

Use either SI (MKS) or CGS as primary units. (SI units are encouraged.) English units may be used as secondary units (in parentheses). An exception would be the use of English units as identifiers in trade, such as  3.5-inch disk drive. 

Avoid combining SI and CGS units, such as current in amperes and magnetic field in oersteds. This often leads to confusion because equations do not balance dimensionally. If you must use mixed units, clearly state the units for each quantity that you use in an equation.

Do not mix complete spellings and abbreviations of units:  Wb/m2  or  webers per square meter,  not  webers/m2.  Spell units when they appear in text:  ...a few henries,  not  ...a few H.

Use a zero before decimal points:  0.25,  not  .25.  Use  cm3,  not  cc. 


\vspace{7pt}
\section{Some Common Mistakes}
The word  data  is plural, not singular.
The subscript for the permeability of vacuum  0, and other common scientific constants, is zero with subscript formatting, not a lowercase letter  o. 
In American English, commas, semi-/colons, periods, question and exclamation marks are located within quotation marks only when a complete thought or name is cited, such as a title or full quotation. When quotation marks are used, instead of a bold or italic typeface, to highlight a word or phrase, punctuation should appear outside of the quotation marks. A parenthetical phrase or statement at the end of a sentence is punctuated outside of the closing parenthesis (like this). (A parenthetical sentence is punctuated within the parentheses.)
A graph within a graph is an  inset,  not an  insert.  The word alternatively is preferred to the word  alternately  (unless you really mean something that alternates).

Do not use the word  essentially  to mean  approximately  or  effectively. 
In your paper title, if the words  that uses  can accurately replace the word using, capitalize the  u ; if not, keep using lower-cased.
Be aware of the different meanings of the homophones  affect  and  effect,   complement  and  compliment,   discreet  and  discrete,   principal  and  principle. 
Do not confuse  imply  and  infer. 
The prefix  non  is not a word; it should be joined to the word it modifies, usually without a hyphen.
There is no period after the  et  in the Latin abbreviation  et al. 
The abbreviation  i.e.  means  that is,  and the abbreviation  e.g.  means  for example. 
An excellent style manual for science writers is \cite{Young}.


\subsection{Subsection Heading Here}
Subsection text here.


\subsubsection{Subsubsection Heading Here}
Subsubsection text here.


% An example of a floating figure using the graphicx package.
% Note that \label must occur AFTER (or within) \caption.
% For figures, \caption should occur after the \includegraphics.
% Note that IEEEtran v1.7 and later has special internal code that
% is designed to preserve the operation of \label within \caption
% even when the captionsoff option is in effect. However, because
% of issues like this, it may be the safest practice to put all your
% \label just after \caption rather than within \caption{}.
%
% Reminder: the "draftcls" or "draftclsnofoot", not "draft", class
% option should be used if it is desired that the figures are to be
% displayed while in draft mode.
%
%\begin{figure}[!t]
%\centering
%\includegraphics[width=2.5in]{myfigure}
% where an .eps filename suffix will be assumed under latex, 
% and a .pdf suffix will be assumed for pdflatex; or what has been declared
% via \DeclareGraphicsExtensions.
%\caption{Simulation Results}
%\label{fig_sim}
%\end{figure}

% Note that IEEE typically puts floats only at the top, even when this
% results in a large percentage of a column being occupied by floats.


% An example of a double column floating figure using two subfigures.
% (The subfig.sty package must be loaded for this to work.)
% The subfigure \label commands are set within each subfloat command, the
% \label for the overall figure must come after \caption.
% \hfil must be used as a separator to get equal spacing.
% The subfigure.sty package works much the same way, except \subfigure is
% used instead of \subfloat.
%
%\begin{figure*}[!t]
%\centerline{\subfloat[Case I]\includegraphics[width=2.5in]{subfigcase1}%
%\label{fig_first_case}}
%\hfil
%\subfloat[Case II]{\includegraphics[width=2.5in]{subfigcase2}%
%\label{fig_second_case}}}
%\caption{Simulation results}
%\label{fig_sim}
%\end{figure*}
%
% Note that often IEEE papers with subfigures do not employ subfigure
% captions (using the optional argument to \subfloat), but instead will
% reference/describe all of them (a), (b), etc., within the main caption.


% An example of a floating table. Note that, for IEEE style tables, the 
% \caption command should come BEFORE the table. Table text will default to
% \footnotesize as IEEE normally uses this smaller font for tables.
% The \label must come after \caption as always.
%
%\begin{table}[!t]
%% increase table row spacing, adjust to taste
%\renewcommand{\arraystretch}{1.3}
% if using array.sty, it might be a good idea to tweak the value of
% \extrarowheight as needed to properly center the text within the cells
%\caption{An Example of a Table}
%\label{table_example}
%\centering
%% Some packages, such as MDW tools, offer better commands for making tables
%% than the plain LaTeX2e tabular which is used here.
%\begin{tabular}{|c||c|}
%\hline
%One & Two\\
%\hline
%Three & Four\\
%\hline
%\end{tabular}
%\end{table}


% Note that IEEE does not put floats in the very first column - or typically
% anywhere on the first page for that matter. Also, in-text middle ("here")
% positioning is not used. Most IEEE journals/conferences use top floats
% exclusively. Note that, LaTeX2e, unlike IEEE journals/conferences, places
% footnotes above bottom floats. This can be corrected via the \fnbelowfloat
% command of the stfloats package.


\vspace{7pt}
\section{Conclusion}
For all further questions please read first the WORD template that is also provided on the EuCAP 2016 website, before consulting the local organising committee.




% conference papers do not normally have an appendix


% use section* for acknowledgement
\vspace{7pt}
\section*{Acknowledgment}


The authors would like to thank...





% trigger a \newpage just before the given reference
% number - used to balance the columns on the last page
% adjust value as needed - may need to be readjusted if
% the document is modified later
%\IEEEtriggeratref{8}
% The "triggered" command can be changed if desired:
%\IEEEtriggercmd{\enlargethispage{-5in}}

% references section

% can use a bibliography generated by BibTeX as a .bbl file
% BibTeX documentation can be easily obtained at:
% http://www.ctan.org/tex-archive/biblio/bibtex/contrib/doc/
% The IEEEtran BibTeX style support page is at:
% http://www.michaelshell.org/tex/ieeetran/bibtex/
%
% argument is your BibTeX string definitions and bibliography database(s)
%
%
% <OR> manually copy in the resultant .bbl file
% set second argument of \begin to the number of references
% (used to reserve space for the reference number labels box)
\vspace{7pt}
\begin{thebibliography}{1}

\bibitem{IEEEhowto:kopka}
H.~Kopka and P.~W. Daly, \emph{A Guide to \LaTeX}, 3rd~ed.\hskip 1em plus
  0.5em minus 0.4em\relax Harlow, England: Addison-Wesley, 1999.

\bibitem{Young}
M.~Young, \emph{The Technical Writer's Handbook}, Mill Valley, CA: University Science, 1989.

\end{thebibliography}




% that's all folks
\end{document}